\documentclass{article}
\usepackage[utf8]{inputenc}
\usepackage{amsmath}
\usepackage{amssymb}
\usepackage{hyperref}
\usepackage{geometry}
\geometry{
 a4paper,
 margin=1in,
}
\begin{document}
\section*{“Nowcasting German Unemployment Using Google Search Data: A MIDAS Approach”}

\begin{itemize}
    \item\textbf{Abstract} : Here you could summarize
    \begin{itemize}
        \item Objective: real-time nowcasting of monthly unemployment.
        \item Method: mixed-frequency MIDAS regressions with weekly Google Trends.
        \item Key result: MIDAS slightly improves forecasts over AR(3), especially in recent periods.
    \end{itemize}

    \item  \textbf{Introduction (1-1.5}: 
    \begin{itemize}
        \item Motivation: Policymakers need early indicators of labor market changes.
        \item Digital data (searches, social media) provide high frequency signals 
        \item Existing literature shows potential but robusteness and bias are debated
        \item Research question: “Can weekly Google search data improve the nowcast of monthly unemployment in Germany compared to an autoregressive benchmark?”
        \item Contribution: Use a recursive r-MIDAS setup with multiple search terms and ridge regularization
    \end{itemize}

    \item  \textbf{Literature review}
    \begin{itemize}
        \item Traditional unemployment forecasting (Phillips curve, ARIMA, etc.).
        \item  High-frequency indicators \& nowcasting (Stock \& Watson, Foroni \& Marcellino, etc.).
        \item Google Trends for labor market prediction (Askitas \& Zimmermann 2009; D’Amuri & Marcucci 2017).
        \item Gaps: need for mixed-frequency frameworks and recursive real-time evaluation.
    \end{itemize}

    \item \textbf{Data (2-3 pages)}:
    \begin{itemize}
        \item Dependent variable: Monthly unemployment rate (Eurostat or Destatis, seasonally adjusted).
        \item High-frequency regressors: Google Trends for 5 keywords — explain selection and scaling issue.
        \item Show your stitching and normalization procedure (include the seam-check table and plot).
        \item Discuss potential bias (scaling windows, overlapping normalization).
        \item Brief stationarity and summary statistics table.
    \end{itemize}

    \item \textbf{Methology: 4-5}
    \begin{itemize}
        \item Baseline: AR(3) on monthly  unemployment.
        \item \begin{equation}
  \Delta u_t \;=\; \alpha
  + \sum_{i=1}^{p} \phi_i \,\Delta u_{t-i}
  + \beta \sum_{k=1}^{K} w_k(a,b)\, x_{t-k}
  + \varepsilon_t .
  \label{eq:midas}
\end{equation}
\item $w_k(a,b)$ = normalized Beta lag polynomial. A common parameterization is
 \begin{equation}
    \tilde w_k(a,b) \;=\;
    \left(\frac{k}{K}\right)^{a-1}\!\left(1-\frac{k}{K}\right)^{b-1},
    \qquad
    w_k(a,b) \;=\; \frac{\tilde w_k(a,b)}{\sum_{j=1}^{K} \tilde w_j(a,b)} .
    \label{eq:beta-weights}
  \end{equation}
\item Discuss ridge penalty across multiple search terms:
\begin{equation}
  \min_{\bm{\beta}}\;\; (y - X\bm{\beta})^{\!\top}(y - X\bm{\beta})
  \;+\; \lambda \lVert \bm{\beta}\rVert_2^2 .
  \label{eq:ridge}
\end{equation}
    \end{itemize}

\item \textbf{Empirical results 4-5}
\begin{itemize}
    \item Stationarity tests (ADF, KPSS).
    \item Stitching sanity checks (show seam jumps, p-values).
    \item Forecast evaluation table (RMSFE, MAE, DM tests).
    \begin{itemize}
        \item Actual vs predicted deltau.
        \item MIDAS weights (Beta curve).
        \item Rolling RMSFE difference (MIDAS – AR3).
    \end{itemize}
    \item Interpretation:
    \begin{itemize}
        \item Small but stable gains from MIDAS.
        \item Improvement concentrated in crisis months.
        \item Ridge stabilizes coefficients without overfitting.
    \end{itemize}
\end{itemize}
\item \textbf{Robustness and overfitting 2-3}
\begin{itemize}
    \item Change info set (rule="eom" vs "mid").
    \item Try different lag lengths (K = 13, 26, 52).
    \item Multi-term vs Factor-MIDAS comparison (brief PCA summary).
    \item Subsample analysis (pre-2020 vs post-2020).
\end{itemize}

\item  \textbf{Conclusion (1 page)}
\begin{itemize}
    \item \textbf{main takeaway}: “High-frequency Google search data modestly improves real-time forecasts of German unemployment.”
    \item Limitations: stationarity of GT series, short history, scaling bias.
    \item Future work: combine MIDAS with ML (LASSO-MIDAS, nonlinear kernels).
\end{itemize}

    
\end{itemize}



\section{Google trend breaks}
\subsection{The main problem}
Google Trends scale breaks across 5-year windows. Each API pull is scaled to its own peak (=100). If you just splice windows, you create artificial jumps at window boundaries → biased levels and spurious dynamics.

\subsection{How did we fix it}
\textbf{Stitching algorithm (window chaining with overlap ratios)}:
The goal was to pull all windows on a common scale by chaining scale factors from overlapping dates.
So what does our switching windows function does is it process windows in a chronological order (by window start, window end, date). For each new window, compute same-date overlaps with the previous window.
And the compute a robust scale factor which is the median of :
\[
\text{factor} = \frac{\text{median}(\text{prev window switched value})}{\text{current window raw value}}over overlaps
\]
Multiply the current window’s raw values by the cumulative scale.
When the same date appears twice (because of overlap), keep the latest window’s stitched value (drop duplicates with keep="last").
Regularize to W-SUN and linearly interpolate tiny gaps.
This works as after switching any two windows that share the same calendar dates must agree in levels (no seam). That’s precisely what we test next.

You can see that as : for every boundary (2014→2015, 2017→2018, 2020→2021, 2023→2024), the mean/median difference = 0.0 on same-date overlaps → no visible seam after stitching. Perfect.

\subsection{Variance explained by windows sanity check}
We also created a window map (win map) aligned to your weekly index and computed the pseudo-ANOVA R² of stitched standardized GT on window dummies:
How could you interprete the 13\%. If switching 
had failed, R² would be very high (windows explain most variance). Your value indicates stitching removed most window-level artifacts while some genuine period differences remain (expected).





\end{document}